\clearpage
\section{Studium wykonalności bezzałogowego samolotu eksperymentalnego}
    \paragraph{}
        W nawiązaniu do metodologii projektowania samolotów eksperymentalnych przedstawionej w literaturze, 
        kolejnym etapem pracy jest przeprowadzenie obliczeń skróconych dla modelu badawczego w zmniejszonej skali. 
        Zgodnie z opisanymi zasadami, w celu ograniczenia kosztów oraz umożliwienia zbadania wybranej konfiguracji aerodynamicznej, 
        zdecydowano się na zaprojektowanie bezzałogowego modelu do badań w locie.
    \paragraph{}
        Do przeliczenia właściwości aerodynamicznych i masowych pomiędzy obiektem w pełnej skali 
        a modelem badawczym zastosowana zostanie teoria podobieństwa dynamicznego, oparta na skalowaniu Froude’a 
        (zgodnie z wytycznymi z tabeli 2.12). 
        Wybór tego kryterium podobieństwa pozwoli na zachowanie odpowiednich relacji sił bezwładności do sił ciężkości, 
        co jest kluczowe podczas badań w locie.
    \paragraph{}
        W dalszej części rozdziału przeprowadzona zostanie analiza doboru masy,
        obciążenia powierzchni oraz napędu dla modelowej skali, 
        z uwzględnieniem ograniczeń wynikających zarówno z budżetu, 
        jak i planowanych parametrów geometrycznych (dążenie do uzyskania jak największej liczby Reynoldsa).